% Gerado a partir de Unidade_1/Resolucoes_e_Pratica/Lista_Treino_Media.md
\chapter{Lista de Treino - Nível Médio (15 Questões)}

\textbf{Disciplina:} Sinais e Sistemas (TEC 412)\\
\textbf{Tópicos:} Aplicações Algébricas, Sinais Discretos e Manipulação de Tempo.

\section{Tópico 1: Sinais Pares e Ímpares}\label{tuxf3pico-1-sinais-pares-e-uxedmpares}

\textbf{Questão 1:} Determine a parte par do sinal \(x(t) = e^{j2t} + \cos(3t)\).

\begin{quote}
\textbf{Resolução:}\\
Expandindo \(e^{j2t} = \cos(2t) + j\sin(2t)\).\\
Então \(x(t) = \cos(2t) + j\sin(2t) + \cos(3t)\).\\
Invertendo: \(x(-t) = \cos(-2t) + j\sin(-2t) + \cos(-3t) = \cos(2t) - j\sin(2t) + \cos(3t)\).\\
Parte Par: \(x_p(t) = \frac{x(t) + x(-t)}{2}\).\\
Ao somar, os senos imaginários se anulam.\\
\(x_p(t) = \frac{2\cos(2t) + 2\cos(3t)}{2} = \cos(2t) + \cos(3t)\).
\end{quote}

\textbf{Questão 2:} O pulso discreto \(x[n] = u[n] - u[n-4]\) não é par. Determine sua parte par \(x_p[n]\).

\begin{quote}
\textbf{Resolução:}\\
O pulso vale 1 para \(n \in \{0, 1, 2, 3\}\).\\
\(x[-n] = u[-n] - u[-n-4]\), que vale 1 para \(n \in \{-3, -2, -1, 0\}\).\\
Parte Par: \(x_p[n] = \frac{x[n] + x[-n]}{2}\).\\
Para \(n=0\): \(x_p[0] = (1+1)/2 = 1\).\\
Para \(n \in \{1, 2, 3\}\) e \(n \in \{-1, -2, -3\}\): apenas uma das parcelas vale 1, logo \(x_p[n] = 0.5\).\\
Para \(n \ge 4\) ou \(n \le -4\): ambas zeram, \(x_p[n]=0\).
\end{quote}

\section{Tópico 2: Energia e Potência}\label{tuxf3pico-2-energia-e-potuxeancia}

\textbf{Questão 3:} Calcule a energia do sinal \(x(t) = e^{-3t}u(t)\).

\begin{quote}
\textbf{Resolução:}\\
\(E = \int_{0}^{\infty} (e^{-3t})^2 dt = \int_{0}^{\infty} e^{-6t} dt\)\\
\(E = \left[ \frac{e^{-6t}}{-6} \right]_0^\infty = 0 - \left( \frac{1}{-6} \right) = \frac{1}{6} \text{ J}\).\\
\textbf{Gabarito: Sinal de Energia (1/6 Joules).}
\end{quote}

\textbf{Questão 4:} Calcule a energia total do sinal discreto \(x[n] = (\frac{1}{2})^n u[n]\).

\begin{quote}
\textbf{Resolução:}\\
Em tempo discreto, Energia é a soma infinita dos quadrados.\\
\(E = \sum_{n=-\infty}^{\infty} |x[n]|^2 = \sum_{n=0}^{\infty} \left( \frac{1}{2} \right)^{2n} = \sum_{n=0}^{\infty} \left( \frac{1}{4} \right)^n\).\\
Essa é uma progressão geométrica de razão \(r = 1/4\).\\
Soma PG infinita: \(S = \frac{1}{1 - r} = \frac{1}{1 - 1/4} = \frac{1}{3/4} = \frac{4}{3}\).\\
\textbf{Gabarito: Energia = 4/3.}
\end{quote}

\section{Tópico 3: Memória}\label{tuxf3pico-3-memuxf3ria}

\textbf{Questão 5:} O integrador genérico \(y(t) = \int_{-\infty}^{t} x(\tau) d\tau\) possui memória?

\begin{quote}
\textbf{Resolução:}\\
A integral computa a área da entrada desde \(-\infty\) até o instante presente \(t\). Isso exige o acúmulo e armazenamento de todos os infinitos valores passados.\\
\textbf{Gabarito: Com Memória.}
\end{quote}

\textbf{Questão 6:} O filtro extrator de pico \(y[n] = \max(x[n], x[n-1])\) é sem memória?

\begin{quote}
\textbf{Resolução:}\\
O sistema precisa comparar o valor atual da entrada (\(x[n]\)) com o valor imediatamente anterior (\(x[n-1]\)). Como requer a leitura de \(n-1\) (passado), o sistema exige registro de armazenamento da última amostra.\\
\textbf{Gabarito: Com Memória.}
\end{quote}

\section{Tópico 4: Causalidade}\label{tuxf3pico-4-causalidade}

\textbf{Questão 7:} O sistema \(y(t) = x(t/3)\) é causal?

\begin{quote}
\textbf{Resolução:}\\
Testando um tempo positivo, digamos \(t=6\).\\
\(y(6) = x(6/3) = x(2)\). Ele precisa do passado. Parece causal.\\
Testando um tempo negativo, digamos \(t=-6\).\\
\(y(-6) = x(-6/3) = x(-2)\).\\
Para calcular o que ocorre hoje no instante \(-6\), eu preciso de uma informação do instante \(-2\). Como \(-2\) é MAIOR que \(-6\), \(-2\) está no futuro relativo!\\
\textbf{Gabarito: NÃO-Causal.}
\end{quote}

\textbf{Questão 8:} O sistema \(y(t) = x(t) \cdot \cos(t+1)\) é causal?

\begin{quote}
\textbf{Resolução:}\\
O ``+1'' está dentro do cosseno, que é o ganho do sistema. Mas não está dentro do ``x(t)''. A entrada exigida para a saída \(y(t_0)\) é APENAS \(x(t_0)\). Não se lê o futuro do sinal de entrada.\\
\textbf{Gabarito: Causal.}
\end{quote}

\section{Tópico 5: Linearidade}\label{tuxf3pico-5-linearidade}

\textbf{Questão 9:} O sistema extrator de parte real \(y(t) = \text{Re}\{x(t)\}\) é linear?

\begin{quote}
\textbf{Resolução:}\\
Seja a combinação complexa genérica \(a \cdot x_1(t) + b \cdot x_2(t)\), onde \(a\) e \(b\) podem ser coeficientes \textbf{complexos}.\\
Se o sistema for linear para qualquer constante complexa, a saída será \(y_3(t) = \text{Re}\{a \cdot x_1(t) + b \cdot x_2(t)\}\).\\
Porém, \(\text{Re}\{a \cdot x\} \neq a \cdot \text{Re}\{x\}\) se `\(a\)' tiver parte imaginária. (Ex: se \(a=j\) e \(x=1\), \(\text{Re}\{j \cdot 1\} = 0\), mas \(j \cdot \text{Re}\{1\} = j\)).\\
A linearidade morre no campo complexo.\\
\textbf{Gabarito: NÃO-Linear.}
\end{quote}

\textbf{Questão 10:} O sistema \(y(t) = \int_{0}^{t} x(\tau) d\tau\) é linear?

\begin{quote}
\textbf{Resolução:}\\
Entrada combinada: \(x_3(t) = ax_1(t) + bx_2(t)\).\\
\(y_3(t) = \int_{0}^{t} [ax_1(\tau) + bx_2(\tau)] d\tau\).\\
Pela propriedade da integral da soma, abrimos em duas:\\
\(y_3(t) = a\int_{0}^{t} x_1(\tau) d\tau + b\int_{0}^{t} x_2(\tau) d\tau\).\\
Que resulta exatamente em \(a y_1(t) + b y_2(t)\).\\
\textbf{Gabarito: Linear.}
\end{quote}

\section{Tópico 6: Invariância no Tempo}\label{tuxf3pico-6-invariuxe2ncia-no-tempo}

\textbf{Questão 11:} O sistema modulador \(y(t) = \sin(t) \cdot x(t)\) é invariante no tempo?

\begin{quote}
\textbf{Resolução:}\\
Atrasando apenas o sinal de entrada: \(y_1(t) = \sin(t) \cdot x(t-t_0)\).\\
Atrasando todas as variáveis temporais do sistema: \(y_2(t) = \sin(t-t_0) \cdot x(t-t_0)\).\\
É claro que \(y_1(t) \neq y_2(t)\). O ganho oscila no tempo.\\
\textbf{Gabarito: Variante no Tempo.}
\end{quote}

\textbf{Questão 12:} O sistema dizimador \(y[n] = x[2n]\) é invariante no tempo?

\begin{quote}
\textbf{Resolução:}\\
1) Atrasa entrada: O sinal \(x_1[n] = x[n-n_0]\) gerará a saída \(y_1[n] = x[2n - n_0]\).\\
2) Atrasa a saída: Substitui \(n\) por \(n-n_0\): \(y_2[n] = y[n-n_0] = x[2(n-n_0)] = x[2n - 2n_0]\).\\
Como \(x[2n - n_0] \neq x[2n - 2n_0]\), ele varia com a época de aplicação.\\
\textbf{Gabarito: Variante no Tempo.}
\end{quote}

\section{Tópico 7: Estabilidade BIBO}\label{tuxf3pico-7-estabilidade-bibo}

\textbf{Questão 13:} O sistema \(y(t) = e^{x(t)}\) é estável?

\begin{quote}
\textbf{Resolução:}\\
Supomos entrada limitada \(-B \le x(t) \le B\).\\
A máxima excursão possível da saída ocorre quando a entrada é \(+B\):\\
\(y(t)_{max} = e^B\).\\
Sendo \(e \approx 2.71\) e \(B\) finito, \(e^B\) é rigorosamente um número finito. Não há como tender ao infinito sem que a própria entrada o faça.\\
\textbf{Gabarito: Estável.}
\end{quote}

\textbf{Questão 14:} O sistema \(y[n] = \log(1 + |x[n]|)\) é estável?

\begin{quote}
\textbf{Resolução:}\\
Como a entrada é restrita em \(|x[n]| \le B\):\\
A máxima saída será \(y_{max} = \log(1 + B)\). O logaritmo cresce muito devagar, e o logaritmo de qualquer número finito positivo sempre resulta num valor finito.\\
\textbf{Gabarito: Estável.}
\end{quote}

\section{Tópico 8: Mistura de Conceitos}\label{tuxf3pico-8-mistura-de-conceitos}

\textbf{Questão 15:} Classifique o sistema com atraso fixo \(y(t) = x(t) + x(t-1)\) quanto aos 4 critérios base.

\begin{quote}
\textbf{Resolução e Gabarito:}

\begin{enumerate}
\def\labelenumi{\arabic{enumi}.}
\tightlist
\item
  Memória: Precisa do instante atual (\(t\)) e do passado (\(t-1\)). \textbf{Com Memória.}
\item
  Causalidade: Como o argumento \(t-1\) nunca será maior que \(t\), ele jamais exigirá o futuro. \textbf{Causal.}
\item
  Linearidade: A soma de sinais combina perfeitamente sem elevar a potências ou inserir constantes isoladas. \textbf{Linear.}
\item
  Invariância: Atrasar a entrada gera \(x(t-t_0) + x(t-t_0-1)\). Atrasar as referências `t' na equação geral gera exatamente \(x(t-t_0) + x((t-t_0)-1)\). \textbf{Invariante no Tempo.}
\end{enumerate}
\end{quote}
